\subsection{Иерархия памяти: скорость, объём и доступ}
\label{lecture04-memory}
Быструю память дорого делать большой. Поэтому компьютер сочетает несколько
уровней: небольшие регистры и кэши, большую RAM и ещё более вместительные
накопители. Эта иерархия --- компромисс, а не закон, что любая операция
на верхнем уровне всегда быстрее любой операции на нижнем.

\begin{center}
\small
\begin{tabular}{|p{0.19\linewidth}|p{0.27\linewidth}|p{0.4\linewidth}|}
\hline
Уровень & Типичная технология & Доступ и назначение \\
\hline
Регистры & Схемы внутри CPU & Указаны в инструкциях; текущие операнды \\
L1/L2/L3 & SRAM & Аппаратно управляемые копии линий памяти \\
RAM & DRAM & Адресуемая рабочая память программ \\
VRAM & Обычно GDDR/HBM & Рабочая память GPU; у встроенного GPU часто общая RAM \\
SSD/HDD & NAND Flash / магнитная запись & Блочный ввод-вывод, долговременное хранение \\
Прошивка & Обычно Flash & Начальная настройка и запуск системы \\
\hline
\end{tabular}
\end{center}

SRAM сохраняет состояние при наличии питания без периодического обновления
заряда, необходимого DRAM; обе технологии энергозависимы. «Статическая»
в названии SRAM не означает сохранение после выключения. Flash энергонезависима.
Регистры не следует механически называть ещё одним уровнем SRAM-кэша.

\paragraph{Ориентиры скорости, а не паспортные константы.}
Для современного настольного компьютера полезны следующие порядки величин:
типичные простые зависимые операции с регистрами и попадания в L1 ---
единицы тактов; L2 --- часто десятки тактов, L3 --- десятки и более;
обращение к DRAM --- десятки--сотни наносекунд; случайный доступ к SSD ---
десятки--сотни микросекунд, к HDD --- миллисекунды. Результат зависит от
процессора, устройства, нагрузки, очередей и попаданий в буферы.
Деление в регистрах, например, не обязано занимать один такт.

При 3 GHz ожидание 100 нс соответствует примерно 300 тактам.
CPU может перекрывать ожидание другой работой, но цепочка зависимых загрузок
часто не позволяет скрыть всю задержку. Доступность также означает
\emph{способ обращения}: инструкция загрузки из обычной памяти не заменяет
команду чтения сектора SSD. Отображённый в память файл может выглядеть для
программы как адресуемые байты, но его подгрузкой управляет ОС.

\subsection{Кэш процессора}
\textbf{Кэш} хранит копии блоков памяти, которые могут понадобиться скоро.
Попадание (cache hit) означает, что нужная линия найдена на проверяемом
уровне; промах (cache miss) требует дальнейшего поиска. Промах L1 не означает
автоматически обращение к RAM: данные могут находиться в L2, L3 или кэше
другого ядра в согласованной системе.

Типичная учебная схема: L1I для инструкций, L1D для данных, собственный L2
ядра и общий для нескольких ядер L3. Эта схема не универсальна. Она
объясняет, почему получение байтов команды и чтение её операнда --- разные
обращения. Программа обычно не распределяет свои переменные по L1/L2/L3
вручную: размещением управляет аппаратура.

\paragraph{Кэш-линия и локальность.}
Для примеров принимаем линию 64 байта --- типичный размер у современных
x86-64. Загрузка одного \texttt{int} может потребовать получения целой линии.
Если \texttt{sizeof(int)==4} и начало группы выровнено по границе линии,
в ней помещаются 16 элементов. Запрос 17 последовательных элементов
затрагивает уже две линии. Если начало смещено, даже 16 элементов могут
занять две линии.

\textbf{Пространственная локальность}: рядом с нужным адресом часто лежат
следующие нужные данные. \textbf{Временная локальность}: недавно использованное
значение может понадобиться снова. При последовательном обходе массива
полезно много байтов каждой линии. При шаге в 16 четырёхбайтовых элементов
каждое чтение попадает в новую линию, и из 64 байтов непосредственно
используются лишь 4. Это иллюстрация использования линий, а не обещание
ускорения ровно в 16 раз: влияют предвыборка, параллельные запросы и объём кэша.

\paragraph{Запись и dirty-линии.}
Для обычной кэшируемой памяти x86-64 распространена политика
\textbf{write-back}. Запись проходит через буфер записей; кэш-линия получает
новое значение и может стать \emph{грязной} (dirty), то есть отличаться от
копии на следующем уровне. Она будет передана дальше по правилам кэша,
например при вытеснении. Следующая загрузка на том же ядре может получить
собственную ещё ожидающую запись из буфера. Это не означает, что программа
произвольно получает старое значение после своего присваивания.

Правила видимости между ядрами сложнее и будут важны при многопоточности.
Главное сейчас: \textbf{запись в переменную, запись в DRAM и сохранение файла
на SSD --- не одно событие}. В нашем примере рассматривается одно ядро,
обычная память и отсутствие конкурирующих изменений.

\subsection{Ввод-вывод, DMA и прерывания}
\label{lecture04-interrupts}
Программа запрашивает действие у ОС. Драйвер настраивает контроллер,
передавая команду, параметры и сведения о буферах. Регистры устройства
могут быть отображены в адресное пространство (MMIO); это \emph{не}
обычные ячейки RAM и не регистры общего назначения CPU. На x86 есть
также отдельные инструкции портового ввода-вывода. Доступ к оборудованию
обычно выполняется привилегированным кодом ОС.

\paragraph{Опрос и уведомление.}
При \textbf{опросе} (polling) CPU периодически проверяет состояние устройства.
Это просто, но частые проверки расходуют процессорное время. При
\textbf{прерывании} (interrupt) система получает уведомление о событии.
Прерывания имеют накладные расходы, поэтому высокопроизводительный
ввод-вывод может сочетать их с опросом и обработкой пакетов событий.

\paragraph{DMA: передача без копирования каждого байта CPU.}
При прямом доступе к памяти (Direct Memory Access) устройство или
контроллер переносит блок в RAM либо из RAM после настройки со стороны
ОС. CPU не обязан выполнять отдельные \texttt{mov} для каждого байта.
Однако он участвует в подготовке операции и обработке результата.
Адрес для устройства не следует отождествлять с пользовательским указателем:
ОС организует допустимый буфер и необходимые преобразования.

Пример чтения, когда нужных данных ещё нет в файловом кэше ОС:
\begin{enumerate}
\item Программа запрашивает чтение; ОС и драйвер готовят запрос и буфер.
\item Контроллер накопителя читает данные и передаёт их в память по DMA.
\item Устройство сообщает о завершении, например прерыванием.
\item ОС отмечает завершение и делает результат доступным программе;
      в обычном буферизованном чтении возможна дополнительная копия.
\end{enumerate}
Если данные уже есть в файловом кэше RAM, физического чтения накопителя
может не потребоваться. Файловый кэш ОС и кэш CPU --- разные уровни.

\paragraph{Как обрабатывается аппаратное прерывание.}
\begin{enumerate}
\item Устройство, контроллер или таймер формирует запрос.
\item Механизм доставки направляет его нужному CPU; маскируемый запрос
      может ожидать, пока процессор сможет его принять.
\item CPU сохраняет предусмотренную архитектурой часть состояния и
      передаёт управление обработчику. Программный код сохраняет
      дополнительные регистры, которые будет использовать.
\item Обработчик выясняет причину, подтверждает или снимает запрос,
      фиксирует результат. Длительную работу можно отложить.
\item При возврате состояние восстанавливается; возможно продолжение
      прерванного кода. ОС может также решить запланировать другой поток.
\end{enumerate}
\textbf{Не каждое прерывание вызывает переключение процесса.} Прерывание
таймера даёт ОС возможность получить управление, учитывать время и принимать
решения планировщика. Прерывание завершения ввода-вывода сообщает о событии,
а не обязательно переносит сами данные.

Для USB-клавиатуры хост-контроллер организует обмен по USB и может уведомить
CPU о результате. Нажатие клавиши не вызывает напрямую функцию
пользовательской программы. Термин USB interrupt transfer тоже не следует
понимать как отдельный провод от клавиши к ядру CPU.

\paragraph{Прерывание, исключение, системный вызов.}
Аппаратное внешнее прерывание асинхронно относительно текущей инструкции.
\textbf{Исключение} связано с её выполнением: например, целочисленное деление
на ноль или страничное исключение. Последнее может быть штатным поводом
подгрузить страницу, а не ошибкой программы. \textbf{Системный вызов} ---
намеренное обращение программы к ОС; на Linux x86-64 обычно используется
инструкция \texttt{syscall}. Это не уведомление от периферии. У разных
архитектур названия и механизмы различаются; здесь важно различать причины.

\subsection{Байты, адреса и объекты C}
В примерах байт содержит 8 бит. \textbf{Адрес} указывает место, а
\textbf{содержимое} --- хранимые там биты. \textbf{Адресное пространство}
--- множество адресов, в котором выполняются обращения. Пользовательская
программа обычно работает с виртуальными адресами, переводимыми аппаратурой
при участии ОС в физические. Наличие 64-битного указателя не означает ни
наличие $2^{64}$ байтов RAM, ни допустимость всех его битовых комбинаций.

Машинное слово --- архитектурно обусловленная единица обработки; термин
нужно употреблять с указанием контекста. В NASM \texttt{word} всегда
16 бит, \texttt{dword} --- 32, \texttt{qword} --- 64. У x86-64 существуют
операции с разными размерами, а не только с 64 битами.

Для обычной Linux x86-64 LP64: \texttt{char} --- 1 байт, \texttt{int} --- 4,
\texttt{long} и указатель --- 8. Это свойства платформы. Например,
в распространённой 64-битной Windows используется LLP64 и \texttt{long}
обычно имеет размер 4 байта. Язык C гарантирует \texttt{sizeof(char)==1},
но число бит в байте определяется \texttt{CHAR\_BIT}.

\begin{lstlisting}[language=C,basicstyle=\ttfamily\small]
int x = 20;
int *p = &x;
*p = 25;
\end{lstlisting}
\texttt{x} --- объект типа \texttt{int}, \texttt{\&x} --- его адрес,
\texttt{p} --- переменная, хранящая указатель, \texttt{*p} --- обращение
к объекту по этому указателю. Присваивание \texttt{*p=25} меняет значение
\texttt{x}, но не адрес \texttt{x}. У самой переменной \texttt{p} также
есть собственное место хранения, если оно требуется реализации.

\paragraph{Массив.}
В \texttt{int a[3]=\{10,20,30\}} элементы идут подряд. Для элемента
с допустимым индексом $i$ байтовый адрес равен
\[
\operatorname{base}+i\cdot\operatorname{sizeof}(\text{элемента}).
\]
При условной базе \texttt{0x1000} и размере 4 байта адреса начала элементов
--- \texttt{0x1000}, \texttt{0x1004}, \texttt{0x1008}; размер массива ---
12 байтов. В выражении C \texttt{a+1} шаг уже масштабируется типом:
не нужно вручную умножать его ещё раз. Массив и указатель --- не один тип:
\texttt{sizeof(a)} здесь равно 12, а размер указателя на нашей платформе --- 8.

\paragraph{Порядок байт.}
Значение \texttt{0x12345678} типа \texttt{uint32\_t} на little-endian x86-64
занимает байты \texttt{78 56 34 12} по возрастающим адресам.
В big-endian было бы \texttt{12 34 56 78}. Меняется порядок байтов,
а не направление записи битов внутри каждого байта.
Для просмотра представления объекта в C можно использовать указатель
на \texttt{unsigned char}. Размеры печатают через \texttt{\%zu}, адреса ---
через \texttt{\%p} с аргументом, приведённым к \texttt{void *}.

\subsection{Сквозной пример: путь инструкций и данных}
\label{lecture04-trace}
Рассмотрим одно изменение массива:
\begin{lstlisting}[language=C,basicstyle=\ttfamily\small]
int a[3] = {10, 20, 30};
a[1] += 5;
\end{lstlisting}
Учебная реализация NASM, если \texttt{RBX} уже содержит адрес массива:
\begin{lstlisting}[basicstyle=\ttfamily\small]
mov eax, [rbx + 4]
add eax, 5
mov [rbx + 4], eax
\end{lstlisting}
Квадратные скобки обозначают обращение к памяти; без них \texttt{RBX}
обозначает регистр. Размер операнда памяти здесь выводится из \texttt{EAX}:
32 бита, то есть 4 байта. Это одна возможная реализация, а не обещание
буквального вывода компилятора. Оптимизатор может выбрать инструкцию
с операндом в памяти или вовсе вычислить известный результат заранее.

\subsubsection*{Шаг 0. От файла к доступным страницам}
Исполняемый файл находится на SSD/HDD. При запуске ОС создаёт отображения
кода и данных в виртуальном адресном пространстве. Весь файл не обязательно
сразу копируется в RAM: страницы могут подгружаться по требованию, а часть
данных создаётся при выполнении. В частности, локальный массив может быть
инициализирован инструкциями после входа в функцию.

Для дальнейших кадров предполагаем: страницы кода и данных уже доступны,
права позволяют исполнение кода и чтение/запись массива, исключений нет.
Условный адрес массива --- \texttt{0x1000}, условный адрес начала фрагмента
кода --- \texttt{0x400000}. Они виртуальные и выбраны только для рисунка.

\subsubsection*{Шаг 1. Процессор получает саму инструкцию}
В архитектурном разборе \texttt{RIP} задаёт адрес очередной инструкции.
Подсистема выборки получает байты кода через L1I. Если линия там отсутствует,
запрос обслуживается нижними уровнями и, при необходимости, RAM.
При холодном обращении удобно рисовать путь
\begin{center}
RAM $\longrightarrow$ нижние уровни кэша $\longrightarrow$ L1I
$\longrightarrow$ декодирование.
\end{center}
Не каждый реальный процессор копирует линию во все уровни именно в таком
порядке. Кэш другого ядра тоже может быть источником. На схеме показываем
смысл иерархии, а не универсальный протокол заполнения.

Для выбранного NASM-кодирования три инструкции имеют байты:
\begin{center}
\small
\begin{tabular}{|l|l|l|}
\hline
Условный адрес & Байты & Инструкция \\
\hline
\texttt{0x400000} & \texttt{8B 43 04} & \texttt{mov eax,[rbx+4]} \\
\texttt{0x400003} & \texttt{83 C0 05} & \texttt{add eax,5} \\
\texttt{0x400006} & \texttt{89 43 04} & \texttt{mov [rbx+4],eax} \\
\hline
\end{tabular}
\end{center}
Все три команды помещаются в одну 64-байтовую линию при данном выравнивании.
В этом примере длина каждой равна 3 байтам, но длина инструкций x86 переменная.
Аппаратура может заранее выбрать несколько команд; это не три обязательных
похода в RAM. Значение \texttt{05} уже содержится в коде инструкции сложения:
отдельная загрузка «переменной 5» из массива не нужна.

\subsubsection*{Шаг 2. Загрузка элемента}
\texttt{mov eax,[rbx+4]} вычисляет адрес \texttt{0x1004}. Перевод адреса
и проверка прав обеспечивают доступ к нужной памяти. Подсистема загрузок
ищет данные через L1D; при промахе --- дальше по иерархии.
В нашей учебной схеме при полном промахе линия данных поступает из RAM
через нижние уровни в L1D, а байты \texttt{14 00 00 00} дают число 20
в \texttt{EAX}. Остальные элементы массива не изменились.

\subsubsection*{Шаг 3. Сложение}
CPU декодирует \texttt{add eax,5}, получает значение регистра и непосредственный
операнд из инструкции. АЛУ вычисляет $20+5=25$. Результат попадает в
\texttt{EAX}; инструкция также обновляет арифметические флаги.
Здесь нет обращения к массиву за операндами сложения: 20 уже в регистре.
До следующей инструкции логическое значение \texttt{a[1]} ещё равно 20.

\subsubsection*{Шаг 4. Запись}
\texttt{mov [rbx+4],eax} направляет 4 байта результата в подсистему памяти.
После разрешённой записи программа должна наблюдать \texttt{a[1]==25}.
В обычной реализации запись проходит через буфер; линия L1D получает
байты \texttt{19 00 00 00}. При write-back копия в DRAM может некоторое
время содержать прежние байты: это не отдельное значение переменной,
которое обычная программа выбирает читать вместо согласованной подсистемы.
Файл на SSD от изменения массива не меняется.

\begin{center}
\small
\begin{tabular}{|l|l|l|}
\hline
Момент & \texttt{EAX} & Логическое содержимое массива \\
\hline
До фрагмента & Не используется & \texttt{10, 20, 30} \\
После загрузки & 20 & \texttt{10, 20, 30} \\
После сложения & 25 & \texttt{10, 20, 30} \\
После записи & 25 & \texttt{10, 25, 30} \\
\hline
\end{tabular}
\end{center}

\paragraph{Проверяемая демонстрация.}
Каталог \texttt{demos/lecture04/} содержит программу C, NASM-функцию
и Docker-окружение. Для проверки выполните из корня репозитория:
\begin{verbatim}
make lecture4-check
\end{verbatim}
Проверяются размеры типов, байты числа, размещение массива и результат
вызова NASM.
В вызываемой функции адрес поступает в \texttt{RDI} по System V AMD64 ABI;
обёртка сохраняет \texttt{RBX}, переносит туда адрес и восстанавливает регистр
перед возвратом. Эти дополнительные команды обеспечивают корректный вызов
из C; три центральные инструкции остаются ровно теми, что разобраны выше.
В демонстрации также проверяется их машинное кодирование.

\subsection{Стек и куча: где живут объекты}
\textbf{Стек} --- область, используемая реализацией для организации вызовов:
адресов возврата, сохранённых регистров, части локальных объектов.
В x86-64 \texttt{RSP} --- указатель стека. \texttt{call} сохраняет адрес
возврата в стеке, \texttt{ret} получает его оттуда. В System V AMD64 многие
аргументы передаются через регистры, поэтому «все аргументы на стеке» неверно.

Локальный объект с автоматическим временем хранения обычно существует
до выхода из блока. Реализация может держать его значение в регистре
или устранить объект: язык C не требует помещать каждую локальную переменную
на стек. Возвращать указатель на уже закончивший жизнь локальный объект
нельзя использовать как способ сохранить его данные.

\textbf{Кучей} обычно называют область и механизмы динамического выделения.
\texttt{malloc} предоставляет блок до явного \texttt{free}; объект не исчезает
только из-за выхода из функции, которая его выделила. Размещение и получение
памяти от ОС зависят от реализации аллокатора.
\begin{lstlisting}[language=C,basicstyle=\ttfamily\small]
int *p = malloc(3 * sizeof *p);
if (p != NULL) {
    p[0] = 10;
    p[1] = 20;
    p[2] = 30;
    p[1] += 5;
    free(p);
    p = NULL;
}
\end{lstlisting}
Нужен \texttt{<stdlib.h>}. Локальная переменная \texttt{p} и выделенный
блок --- разные объекты с разным временем жизни. После \texttt{free}
разыменование старого указателя недопустимо; присваивание \texttt{NULL}
этой переменной не исправляет другие копии указателя.
Глобальные объекты и статические локальные объекты имеют статическое время
хранения: деление «всё либо стек, либо куча» неполно.

\subsection{Компьютер целиком: чтение, вычисление, вывод}
Представим программу, которая читает массив чисел, изменяет их и рисует
диаграмму. Ввод-вывод загружает данные с накопителя в память; CPU получает
инструкции и данные через кэши. При подходящей большой задаче часть работы
можно передать GPU. Результат для отображения попадает в графическую
подсистему; для сохранения файла программа отдельно запрашивает запись у ОС.

Узкое место зависит от этапа: случайные чтения HDD ограничены механикой,
большая передача --- пропускной способностью, зависимые чтения RAM ---
задержкой, вычислительный цикл --- исполнительными ресурсами. Переход на SSD
не ускорит непосредственно сложение уже находящихся в регистрах операндов.
Больше RAM помогает при нехватке рабочей памяти; больше VRAM не заменяет
вычислительные блоки GPU. Оценивать нужно путь конкретной задачи.

\subsection*{Итоги и самопроверка}
\begin{enumerate}
\item Откуда берутся байты \texttt{mov eax,[rbx+4]}, а откуда загружаемое число?
\item Почему промах L1 не обязательно приводит к чтению DRAM?
\item Чем отличаются адрес \texttt{0x1004} и число 20?
\item Почему изменённый массив не сохраняется автоматически в файл?
\item Обязательно ли CPU копирует каждый байт при DMA?
\item Обязательно ли прерывание переключает процесс?
\item Почему общая память команд и данных не означает исполняемость стека?
\item Где находится объект, на который указывает локальная переменная \texttt{p}?
\end{enumerate}
\textbf{Краткие ответы:} код выбирается через подсистему инструкций, число ---
через подсистему данных; существуют нижние уровни кэша; адрес указывает место,
число является значением; нужна отдельная операция ввода-вывода; нет,
передачу выполняет устройство/контроллер; нет, это решение ОС; действуют
права страниц; зависит от того, на что указывает \texttt{p}, а не от
места хранения самого указателя.

\newpage
\subsection{Дополнительно: детали за пределами 90 минут}
\paragraph{Линия, индекс и тег.}
Для 64-байтовой линии младшие 6 бит байтового адреса задают смещение внутри
линии. В учебном кэше с прямым отображением и 64 наборами следующие 6 бит
выбирают набор, остальные образуют тег. Тег позволяет отличить разные
блоки, претендующие на одно место. Например, для адреса \texttt{0x1234}
смещение равно \texttt{0x34}=52, номер набора --- 8, тег --- 1.
Это модель кэша, а не описание конкретного L1 современного x86.

\paragraph{Ассоциативность и вытеснение.}
В наборно-ассоциативном кэше один набор может хранить несколько
линий с одним индексом. Английский термин --- set-associative cache.
Такая организация уменьшает некоторые конфликтные промахи.
При нехватке места выбирается линия для вытеснения. Политика не обязана
быть точным LRU. Три причины промаха: первое обращение к линии,
нехватка общего объёма и конкуренция за набор.

\paragraph{Политики записи.}
Write-through передаёт запись также на следующий уровень, write-back
позволяет отложить её до записи изменённой линии. Write-allocate описывает
получение линии при промахе записи; это другой вопрос. Обычная память,
MMIO и специальные инструкции могут иметь разные правила кэширования.

\paragraph{Согласованность кэшей.}
Несколько ядер могут иметь копии одной линии. Протокол согласованности
помогает организовать владение и распространение изменений. Даже изменения
разных переменных в одной линии могут вызывать лишние передачи
(false sharing). Согласованность кэшей не заменяет атомарные операции,
синхронизацию и модель памяти языка C.

\paragraph{Страницы и TLB.}
Виртуальное адресное пространство отображается на физическую память
страницами. В примере со страницей 4 KiB младшие 12 бит задают смещение,
старшие --- номер виртуальной страницы. Для \texttt{0x1234} это страница 1
и смещение \texttt{0x234}=564. TLB кэширует переводы адресов и сведения о
доступе, а не данные массива. Промах TLB может потребовать обхода таблиц;
страничное исключение возникает по иной причине, например при отсутствии
допустимого отображения. Размер страницы зависит от режима системы.

\paragraph{Выравнивание и структуры.}
У структуры \texttt{struct S \{ char c; int x; \};} при обычных правилах
выравнивания Linux x86-64 поле \texttt{x} имеет смещение 4, размер структуры
--- 8 байтов. Между полями есть padding. Для конкретной реализации это
проверяют через \texttt{offsetof} и \texttt{sizeof}; это не универсальный
результат для всех компиляторов и настроек упаковки.

\paragraph{Оптимизация и наблюдаемое поведение.}
Компилятор обязан сохранять требуемое языком наблюдаемое поведение, а не
по одной инструкции на каждую строку. Если результат массива нигде не
используется, его и вычисление могут исчезнуть. Код демонстрации вызывает
внешнюю NASM-функцию и проверяет результат; это делает разбор воспроизводимым.
\texttt{volatile} не следует использовать как универсальное средство
управления кэшем или синхронизации потоков.

\paragraph{Доставка событий.}
На современных x86 запросы могут доставляться через APIC, в том числе
сообщениями MSI/MSI-X от PCIe-устройств. ОС настраивает маршрутизацию и
обработчики. Эти детали не меняют базового различия между передачей данных
и уведомлением о готовности. Реальные системы используют очереди, пакетную
обработку и отложенную работу, чтобы уменьшать накладные расходы.

\paragraph{Источники.}
Общая модель компьютера и иерархия памяти: CS:APP~\cite{bryant-ohallaron-csapp}
и Patterson--Hennessy~\cite{patterson-hennessy}. Регистры, инструкции,
страницы и прерывания: Intel SDM~\cite{intel-manual}. Синтаксис и размеры
операндов: документация NASM~\cite{nasm-doc}. Объекты C, массивы и время
хранения: N1570~\cite{c11-draft}; соглашение вызова демонстрации:
System V AMD64 psABI~\cite{amd64-psabi}.
